\section{Context and claim boundary}

\paragraph{H\'enon dynamics and scaling.}
The source H\'enon model motivates an area-preserving quadratic map and its
cubic generating function \citep{Wang2026HenonModel}. The immediate input to this paper is the exact
Poisson boundary identity for its adelic chirp. Under dilation,
\begin{equation}\label{eq:dilation}
D_aM_{P_6}D_a^{-1}=M_{P_a},
\qquad P_a(x)=2a^3x^3-ax=P_6(ax).
\end{equation}
The coefficient functionals indexed by $a$ span an infinite-dimensional
space before the Poisson map. A finite static projection-rank bound cannot
therefore be promoted to a finite-channel dynamical theorem. Equation
\eqref{eq:dilation} instead points to Mellin transformation in $a$.

\paragraph{Mellin and Poisson analysis.}
Mellin continuation, contour rotation, and generalized hypergeometric
functions are classical tools; we follow standard conventions summarized by
the NIST Digital Library of Mathematical Functions
\citep{DLMFMellin,DLMFHyper}. The adelic Poisson/scaling framework and its
spectral interpretation originate in the Connes program
\citep{Connes1999,ConnesConsani2020}. We do not claim those mechanisms as
new. Our question is narrower: what divisor is forced when the H\'enon cubic
boundary family is Mellin diagonalized without fitted normalization?

\paragraph{Certified complex analysis.}
The proof uses Arb midpoint-radius arithmetic through the python-flint
interface. Arb provides rigorous enclosures rather than ordinary
floating-point estimates \citep{Johansson2017Arb}. The software evaluation
is only one part of the proof: a symbolic contour formula and an analytic
second-derivative majorant convert the enclosures into a Rouch\'e count.

\paragraph{Scope.}
We call $S_H$ a \emph{formal scattering symbol}. No wave operators,
trace-class resolvent difference, or Hilbert-space scattering pair is
constructed here. The word ``determinant'' refers to the pointwise
$2\times2$ matrix determinant unless explicitly qualified. These
distinctions prevent a symbol-level functional equation from being mistaken
for a Hilbert--P\'olya operator theorem.
